% СХЕМА 1 — «Картина мира»: солист и виртуальный оркестр
\begin{tikzpicture}[
  every node/.style={font=\small},
  block/.style={draw, thick, minimum width=2.6cm, minimum height=1.4cm,
                align=center, rounded corners=2pt, fill=white},
]
  % Solist
  \node[block] (solo) at (0,0) {%
    \textbf{Солист}\\
    \footnotesize(пианино)};
  \node at ($(solo.north) + (0,0.5)$) {\footnotesize\itshape играет в своём темпе};

  % Orchestra
  \node[block] (orch) at (6,0) {%
    \textbf{Виртуальный}\\
    \textbf{оркестр}\\
    \footnotesize(колонки)};
  \node at ($(orch.north) + (0,0.5)$) {\footnotesize\itshape подстраивается};

  % Arrows
  \draw[-{Latex[length=3mm]}, thick] (solo.east) -- (orch.west)
        node[midway, above] {\scriptsize MIDI / audio};
  \draw[-{Latex[length=3mm]}, thick, dashed] (orch.south)
        to[bend right=40] node[midway, below=2pt] {\scriptsize звучит синхронно} (solo.south);
\end{tikzpicture}
